% Claim statement file for row claim_3_12 -- "HardInterventionNodeSplitOrder".
% Auto-generated stub by scaffold/solving_current_row.py at row start. The
% manager agent (or a worker it dispatches) fills the body below with the
% claim's statement(s) from the lecture notes -- statement only, no proof
% (proofs live in the sibling claim_*_proof_*.tex file).
%
% Compiles standalone via the `subfiles` package.

\documentclass[main]{subfiles}

\begin{document}
% ref: claim_3_12 (statement)
% title: HardInterventionNodeSplitOrder

\begin{Rem}
    Note that if $W_1$ and $W_2$ are not disjoint and $w \in W_1 \cap W_2 \ins V$
    then first hard intervening on $w$ turns $w$ into an input node, for now indicated as $w^i$, and a node-splitting hard intervention (if we would define it for input nodes) would not change $w^i$.
    If, on the other hand, we would first split the node $w$ into $w^o$ and $w^i$ then we would first need to resolve the ambiguity on which of those two the hard intervention should be applied. A hard intervention on $w^i$ would not do anything, but would leave the additional output node $w^o$ in the graph, while hard intervening on $w^o$ would turn $w^o$ into an additional input node, for now indicated as $(w^o)^i$. So in the latter case we are left with two input node $(w^o)^i$, which does not have any edges, and $w^i$, which might have outgoing edges.
\end{Rem}

\end{document}
